\chapter{Geometría a partir de la Estructura Lineal}

En el Capítulo I, los espacios vectoriales se introdujeron como objetos puramente algebraicos.
Un vector era un elemento abstracto de un conjunto dotado de suma y multiplicación por escalares,
y las transformaciones lineales eran funciones que preservaban dicha estructura.
No se asumió ni se requirió ninguna noción de longitud, ángulo, distancia u ortogonalidad.

\medskip

El álgebra por sí sola nos permite decidir si los vectores son independientes,
si un conjunto genera un espacio,
o si una transformación lineal es inyectiva o sobreyectiva.
Sin embargo, no permite formular preguntas geométricas.
En muchas situaciones, no basta con saber \emph{qué} vectores son posibles;
también necesitamos saber cuán \emph{grandes} son, cuán \emph{próximos} están entre sí,
o cuán \emph{alineados} se encuentran.
Estas nociones se expresan mediante longitud, distancia y ángulo.
Sin estructura adicional, el lenguaje de los espacios vectoriales no distingue
entre «mayor» y «menor», ni permite hablar de perpendicularidad o de la
«mejor aproximación» de un vector mediante otros
(como en las proyecciones y los mínimos cuadrados).
No existe ninguna manera intrínseca, dentro de un espacio vectorial desnudo,
de comparar el tamaño de dos vectores o de determinar si apuntan en direcciones
perpendiculares.

\medskip

La geometría comienza únicamente cuando se impone estructura adicional.

\medskip

Conviene recordar ahora cómo se representa habitualmente el espacio geométrico
en coordenadas.
Dados conjuntos $A_1,\dots,A_n$, su \emph{producto cartesiano}
(también llamado \emph{producto de conjuntos}) es
\[
A_1\times\cdots\times A_n = \{(x_1,\dots,x_n)\mid x_i\in A_i\}.
\]
En particular, el espacio coordenado $\mathbb{R}^n$ es el producto cartesiano de
$n$ copias de $\mathbb{R}$.
Esta construcción por sí sola no es todavía geometría: proporciona un
\emph{modelo} conveniente en el que los «puntos» se representan mediante
$n$-uplas de números reales.
La geometría emerge únicamente después de especificar una regla adicional que
nos permita medir y comparar dichas $n$-uplas.

\medskip

Un tema central de este documento es que la geometría no reemplaza a la
estructura lineal, sino que la enriquece.
Los espacios vectoriales introducidos anteriormente permanecen inalterados;
lo que cambia es que se equipan ahora con una regla que permite a los vectores
interactuar de un modo que produce cantidades escalares.

\medskip

Informalmente, esta regla toma dos vectores y devuelve un escalar.
A partir de esta única operación emergen de manera natural nociones como
longitud, ángulo y distancia.

\medskip

Para precisar esta interacción, introducimos la noción de aplicación bilineal.

\medskip

Sean $E$ y $G$ espacios vectoriales sobre el mismo cuerpo $K$.
Una aplicación
\[
B : E \times G \longrightarrow K
\]
se llama \emph{bilineal} si es lineal en cada argumento por separado.
Es decir, para $y \in G$ fijo, la aplicación
\[
x \longmapsto B(x,y)
\]
es una forma lineal sobre $E$, y para $x \in E$ fijo, la aplicación
\[
y \longmapsto B(x,y)
\]
es una forma lineal sobre $G$.

\medskip

Explícitamente, para todo $x_1,x_2 \in E$, $y_1,y_2 \in G$, y todo escalar
$\alpha \in K$,
\[
B(x_1 + x_2, y) = B(x_1, y) + B(x_2, y), \qquad
B(\alpha x, y) = \alpha B(x, y),
\]
\[
B(x, y_1 + y_2) = B(x, y_1) + B(x, y_2), \qquad
B(x, \alpha y) = \alpha B(x, y).
\]

\medskip

Cuando $E = G$, dicha aplicación se llama \emph{forma bilineal} sobre $E$.

\medskip

A este nivel no se asume ninguna interpretación geométrica.
Una forma bilineal es simplemente una regla que asigna un escalar a un par
ordenado de vectores, sujeta a la linealidad en cada argumento.
No tiene por qué ser simétrica, positiva, ni estar relacionada con ninguna
noción de longitud o ángulo.

\medskip

No obstante, las formas bilineales desempeñan un papel estructural fundamental.
Proporcionan una manera sistemática de asociar vectores con formas lineales
y de construir cantidades escalares a partir de pares de vectores.

\medskip

Las nociones geométricas familiares emergen únicamente después de imponer
condiciones adicionales sobre la forma bilineal.
Estos casos especiales se introducirán en breve.

% CORRECCIÓN 3: La restricción a R se introduce ahora aquí, inmediatamente después
% de las formas bilineales y antes de cualquier mención a positividad o productos
% internos. Anteriormente aparecía mucho más tarde, después de que el producto
% interno ya se había usado implícitamente.

\medskip

Hasta este punto, el cuerpo escalar $K$ ha sido arbitrario.
La estructura lineal pura no requiere la capacidad de comparar escalares.
La geometría, sin embargo, depende de la positividad: expresiones como
$\langle x,x\rangle > 0$ y desigualdades de la forma $r>0$
presuponen una relación de orden.
Por esta razón, restringimos ahora la atención al sistema de los números reales
\[
\Big\langle \mathbb{R}, +, \cdot, \le \Big\rangle,
\]
visto como cuerpo ordenado.
Es precisamente el orden en $\mathbb{R}$ el que permite que las longitudes,
las distancias y la ortogonalidad adquieran significado geométrico.

\medskip

Especializamos ahora las formas bilineales introducidas anteriormente.

\medskip

Sea $E$ un espacio vectorial sobre $\mathbb{R}$.
Una \emph{forma bilineal} sobre $E$ es una función
\[
B : E \times E \longrightarrow \mathbb{R}
\]
que es lineal en cada argumento por separado.
Es decir, para todo $x,x',y,y' \in E$ y todo $\lambda \in \mathbb{R}$,
\[
B(x+x',y)=B(x,y)+B(x',y), \qquad
B(\lambda x,y)=\lambda B(x,y),
\]
y análogamente en el segundo argumento.

\medskip

Las formas bilineales surgen de manera natural a partir de las transformaciones lineales.
Si
\[
\psi : E \longrightarrow E^*
\]
es una aplicación lineal, entonces la fórmula
\[
B(x,y) = \psi(x)(y)
\]
define una forma bilineal sobre $E$.
Recíprocamente, toda forma bilineal determina tal aplicación.
Por tanto, las formas bilineales codifican una manera controlada en que los
vectores pueden interactuar para producir escalares.

\medskip

A este nivel no está presente ninguna geometría.
Una forma bilineal puede ser degenerada, asimétrica o indefinida.
Permite que los vectores interactúen, pero todavía no mide nada.

\medskip

La geometría aparece cuando la forma bilineal satisface restricciones adicionales.

% CORRECCIÓN 1 y 2: El producto interno se define ahora antes de cualquier
% expresión coordenada que lo involucre. La sección sobre la transpuesta, que
% anteriormente interrumpía el camino de las formas bilineales al producto interno,
% ha sido reubicada después de que el producto interno esté establecido. La
% expresión coordenada del producto interno sigue a la definición, no la precede.

\medskip

Un \emph{producto interno} sobre un espacio vectorial $E$ sobre $\mathbb{R}$
es una forma bilineal
\[
\langle \cdot,\cdot \rangle : E \times E \longrightarrow \mathbb{R}
\]
que es simétrica, positiva y no degenerada.
Estas propiedades garantizan que
\[
\langle x,x \rangle > 0 \quad \text{para todo } x \neq 0.
\]

\medskip

Una vez fijada dicha regla, los vectores adquieren propiedades medibles.
Dos vectores pueden declararse ortogonales.
A un vector se le puede asignar una longitud.
La distancia entre dos puntos puede definirse en términos de los vectores que
los unen.
Ninguno de estos conceptos existe antes de la introducción de esta estructura adicional.

\medskip

Como en el Capítulo I, las coordenadas no desempeñan ningún papel fundamental.
Las cantidades geométricas deben ser independientes de cualquier sistema de
referencia particular.
Una fórmula que cambia al cambiar la base describe una representación,
no un objeto geométrico.

\medskip

Solo después de elegir una base aparecen las expresiones coordenadas familiares.
Las sumas de cuadrados, los productos escalares y las ecuaciones cuadráticas
surgen como manifestaciones a nivel coordenado de la estructura subyacente,
no como definiciones de ella.

\medskip

\textit{Expresión coordenada del producto interno.}

Sea $(E,\langle\cdot,\cdot\rangle)$ un espacio con producto interno de dimensión
finita y sea $\{e_1,\dots,e_n\}$ una base ortonormal.
Todo vector $x\in E$ puede escribirse de manera única como
\[
x=\sum_{i=1}^n x_i e_i,
\qquad
y=\sum_{i=1}^n y_i e_i.
\]

Usando la bilinealidad y la condición de ortonormalidad
$\langle e_i,e_j\rangle=\delta_{ij}$, obtenemos
\[
\langle x,y\rangle
=
\sum_{i=1}^n x_i y_i.
\]

Así, en una base ortonormal el producto interno abstracto queda representado
por el familiar \emph{producto escalar} de vectores coordenados en $\mathbb{R}^n$.

\medskip

Esta observación explica el papel del producto de matrices.
Si una transformación lineal $\psi:E\to G$ está representada por una matriz
$A=(a_{ij})$ y un vector $x$ tiene vector coordenado $x=(x_1,\dots,x_n)^T$,
entonces la $i$-ésima componente del producto $Ax$ es
\[
(Ax)_i=\sum_{j=1}^n a_{ij}x_j.
\]

Cada componente es por tanto el producto escalar de la $i$-ésima fila de $A$
con el vector $x$.
El producto de matrices puede así entenderse como la evaluación sistemática de
productos internos entre filas de una matriz y columnas de otra.

\medskip

En este sentido, la regla algebraica familiar para multiplicar matrices no es
una convención computacional arbitraria: admite una interpretación geométrica
natural como expresión coordenada de interacciones medidas por un producto interno.

\medskip

Este punto de vista explica por qué objetos como circunferencias, esferas y
elipsoides pueden describirse mediante ecuaciones algebraicas y, sin embargo,
conservan un significado geométrico invariante bajo cambios de coordenadas.
Las ecuaciones dependen de la base elegida; la geometría no.

\medskip

El objetivo es por tanto mostrar cómo la geometría emerge de la estructura
lineal una vez que se especifica una interacción adecuada entre los vectores.

\medskip

% CORRECCIÓN 2 (continuación): La sección sobre la transpuesta y la transformación
% dual se ubica ahora aquí, después de definir el producto interno. Esto preserva
% la derivación algebraica de la transpuesta como intrínseca (sin necesidad de
% producto interno), a la vez que permite que la observación sobre la dualidad
% interna aparezca de manera natural en el contexto de un producto interno ya establecido.

Aunque la geometría ha comenzado a emerger a través de la interacción bilineal,
el mecanismo algebraico subyacente sigue siendo la transformación lineal.
Para comprender cómo interactúa la estructura geométrica con las aplicaciones
lineales, examinamos una construcción asociada de manera natural a toda
transformación lineal.

\medskip

El objeto central de este trabajo es la transformación lineal. Antes de
introducir nociones geométricas como longitud, ángulo y distancia, hacemos
explícita una construcción fundamental asociada a toda transformación lineal:
su transpuesta.

\medskip

Sean $E$ y $G$ espacios vectoriales sobre el mismo cuerpo $K$, y sea
\[
\psi : E \longrightarrow G
\]
una transformación lineal. Recordemos que los espacios duales $E^*$ y $G^*$
consisten en todas las formas lineales sobre $E$ y $G$, respectivamente.

\medskip

Para cada $g^* \in G^*$, la composición
\[
g^* \circ \psi : E \longrightarrow K
\]
es una forma lineal sobre $E$. Esto define una aplicación
\[
\psi^* : G^* \longrightarrow E^*,
\qquad
\psi^*(g^*) = g^* \circ \psi.
\]

\medskip

La aplicación $\psi^*$ es lineal y se llama la \emph{transpuesta}
(o transformación dual) de $\psi$. La dirección de la flecha se invierte:
\[
E \xrightarrow{\;\psi\;} G
\quad\Longrightarrow\quad
G^* \xrightarrow{\;\psi^*\;} E^*.
\]

\medskip

Esta construcción es intrínseca: depende únicamente de la transformación
lineal $\psi$ y no requiere ninguna elección de base, coordenadas o producto interno.

\medskip

Cuando se eligen bases para $E$ y $G$, y se eligen las bases duales
correspondientes para $E^*$ y $G^*$, la transformación lineal $\psi$ queda
representada por una matriz $A$, y la transformación dual $\psi^*$ queda
representada por la matriz transpuesta $A^T$.
Así, la transposición matricial familiar es la expresión coordenada de la
transformación dual intrínseca $\psi^*$.

\medskip

El producto interno hace más que asignar longitudes.
Establece una identificación canónica entre el espacio vectorial y su dual.

\medskip

Para cada vector $x \in E$, definimos la aplicación
\[
\Phi(x) : E \longrightarrow \mathbb{R},
\qquad
\Phi(x)(y) = \langle x,y \rangle.
\]

Por bilinealidad, $\Phi(x)$ es una forma lineal sobre $E$.
Obtenemos así una aplicación
\[
\Phi : E \longrightarrow E^*.
\]

\medskip

Si el producto interno es no degenerado, esta aplicación es inyectiva.
En dimensión finita, como $\dim E = \dim E^*$, es por tanto un isomorfismo.

\medskip

Así, una vez fijado un producto interno, cada vector determina una forma lineal
única, y toda forma lineal proviene de un vector único.
La distinción entre $E$ y $E^*$ desaparece a nivel estructural.

\medskip

Es precisamente esta identificación la que permite considerar la transpuesta
de una transformación lineal como un operador sobre el propio $E$.
Cuando se introduce la geometría, la dualidad se vuelve interna.

\medskip

Una vez fijado un producto interno, los vectores adquieren longitud.
La norma de un vector $x \in E$ se define por
\[
\|x\| = \sqrt{\langle x,x \rangle}.
\]

\medskip

Así, el espacio vectorial $E$ introducido en el Capítulo I queda ahora dotado
de estructura geométrica adicional:
\[
\Big\langle
\langle \mathbb{R},+,\cdot\rangle,\;
\langle E,+\rangle,\;
\circ,\;
\langle\cdot,\cdot\rangle
\Big\rangle.
\]

El producto interno no es necesario para la linealidad en sí, pero permite
introducir nociones geométricas como longitud, ángulo, ortogonalidad y
transformaciones adjuntas.

\medskip

\textit{Expansión algebraica de la norma.}

Para cualesquiera vectores $x,y \in E$, el cuadrado de la norma de su suma satisface
\[
\|x+y\|^2
=
\langle x+y, x+y \rangle.
\]
Por bilinealidad y simetría del producto interno,
\[
\langle x+y, x+y \rangle
=
\langle x,x \rangle
+
\langle x,y \rangle
+
\langle y,x \rangle
+
\langle y,y \rangle
=
\|x\|^2
+
2\langle x,y \rangle
+
\|y\|^2.
\]

Por tanto,
\[
\|x+y\|^2
=
\|x\|^2
+
\|y\|^2
+
2\langle x,y \rangle.
\]

Esta identidad se cumple en todo espacio con producto interno.
Expresa cómo el término de interacción $\langle x,y \rangle$ mide la
desviación respecto a la aditividad pura de los cuadrados de las longitudes.

\medskip

El producto interno satisface también la desigualdad
\[
|\langle x,y \rangle|
\le
\|x\|\,\|y\|,
\]
conocida como la \emph{desigualdad de Cauchy--Schwarz}.
Como consecuencia,
\[
\|x+y\|^2
\le
\|x\|^2
+
\|y\|^2
+
2\|x\|\,\|y\|
=
(\|x\|+\|y\|)^2.
\]

Tomando raíces cuadradas se obtiene la \emph{desigualdad triangular}
\[
\|x+y\|
\le
\|x\|
+
\|y\|.
\]

\medskip

La \textit{identidad de Pitágoras} surge como el caso especial
$\langle x,y \rangle = 0$,
en el que el término de interacción se anula y
\[
\|x+y\|^2
=
\|x\|^2
+
\|y\|^2.
\]

\medskip

La ortogonalidad también adquiere significado.
Se dice que dos vectores $x,y \in E$ son ortogonales si
\[
\langle x,y \rangle = 0.
\]

\medskip

El orden de los vectores es irrelevante.
La simetría del producto interno implica
\[
\langle x,y \rangle = \langle y,x \rangle,
\]
de modo que la ortogonalidad es una relación mutua.

\medskip

El vector cero es ortogonal a todo vector de $E$.
En efecto, para cualquier $v \in E$,
\[
\langle 0,v \rangle
=
\langle 0\cdot v, v \rangle
=
0\langle v,v \rangle
=
0.
\]
Más aún, el vector cero es el único con esta propiedad.

\medskip

Así, la ortogonalidad es precisamente la condición bajo la cual el término
de interacción se anula.
En ese caso,
\[
\|x+y\|^2 = \|x\|^2 + \|y\|^2.
\]
Los cuadrados de las longitudes se vuelven aditivos exactamente cuando los
vectores no interactúan.

\medskip

% CORRECCIÓN 4: El interludio sobre vectores propios se ha trasladado a aquí,
% después de establecer la norma, la ortogonalidad y el adjunto. Anteriormente
% aparecía entre la identidad de Pitágoras y la sección sobre distancia,
% interrumpiendo la cadena norma -> distancia -> topología en su punto más crítico.

\textit{Interludio: Direcciones invariantes}

Hasta este punto, las transformaciones lineales se han tratado algebraicamente:
una aplicación $\psi : E \to E$ envía vectores a vectores preservando la
estructura lineal.
Todavía no se ha atribuido ningún significado geométrico a esta acción.

\medskip

Una vez introducido un producto interno, sin embargo, la acción de una
transformación lineal puede examinarse geométricamente.
Dado un vector $x \in E$, la imagen $\psi(x)$ puede diferir de $x$ de dos
maneras independientes: su longitud puede cambiar y su dirección puede cambiar.

\medskip

Para un vector genérico, ambos efectos ocurren simultáneamente.
La mayoría de los vectores son estirados y rotados.

\medskip

Se produce un fenómeno notable cuando existe un vector no nulo $x \in E$
tal que
\[
\psi(x) = \lambda x
\]
para algún escalar $\lambda$.
En este caso, la transformación actúa sobre $x$ mediante un reescalado puro.
La dirección de $x$ se preserva.

\medskip

Tal vector se llama \emph{vector propio} de $\psi$, y el escalar $\lambda$
es su \emph{valor propio} correspondiente.
Los vectores propios identifican direcciones en el espacio que son
geométricamente invariantes bajo la transformación.

\medskip

Esta invariancia es una propiedad intrínseca de la propia transformación lineal.
Cambiar la base puede alterar la representación de $\psi$, pero no puede
destruir ni crear vectores propios.

\medskip

Bajo la identificación $E \cong E^*$ inducida por el producto interno, la
transpuesta $\psi^*$ corresponde a un operador lineal sobre $E$, llamado
el \emph{adjunto} de $\psi$.

\medskip

Cuando una transformación lineal proviene de una forma bilineal simétrica (o,
equivalentemente, de un operador autoadjunto), sus vectores propios
correspondientes a valores propios distintos son ortogonales.
Estas direcciones definen un sistema de coordenadas preferido dictado por la
geometría, no por una elección arbitraria.

\medskip

Es en este sentido que los vectores propios sirven como ejes principales de
los elipsoides asociados a las formas cuadráticas.
La diagonalización es el acto de alinear las coordenadas con las direcciones
invariantes ya presentes en la estructura.

\medskip

% CORRECCIÓN 4 (continuación): Con los vectores propios establecidos aquí,
% la sección sobre distancia y topología sigue sin interrupción.

\textit{Distancia y emergencia de la geometría.}

Una vez disponible un producto interno sobre un espacio vectorial $E$, este
determina una norma, y a partir de la norma obtenemos una noción de distancia
entre vectores,
\[
d(x,y)=\|x-y\|.
\]

En este punto no se introduce ninguna estructura nueva: la distancia es
simplemente una reformulación de datos algebraicos ya presentes.

\medskip

Esta distancia nos permite describir la proximidad. Para un vector $x \in E$
y un número real $r > 0$, consideramos el conjunto
\[
B(x,r) = \{\, y \in E \mid d(x,y) < r \,\}.
\]
Dicho conjunto se llama \textit{bola abierta} centrada en $x$ con radio $r$.
Geométricamente, consiste en todos los vectores que se encuentran más cerca
de $x$ que el umbral prescrito.

\medskip

La desigualdad estricta es esencial. Garantiza que todo punto interior a la
bola admite puntos cercanos adicionales, de modo que no se incluye ninguna
frontera.
Este «margen de holgura» es lo que en última instancia nos permite hablar de
continuidad y deformación.

\medskip

La colección de todas las bolas abiertas determina qué subconjuntos de $E$
deben llamarse abiertos: un conjunto es abierto si, alrededor de cada uno de
sus puntos, contiene alguna bola abierta.
De este modo, la distancia genera una noción de apertura sin axiomas adicionales.

% CORRECCIÓN 5: Una frase de transición prepara ahora al lector para la
% definición axiomática de espacio topológico, conectándola con la intuición
% geométrica construida justo antes, en lugar de introducir los axiomas
% abruptamente.

\medskip

Esta noción de apertura puede capturarse mediante un pequeño conjunto de
axiomas que preservan sus propiedades esenciales y permiten extender el concepto
mucho más allá del contexto métrico.
La estructura resultante es lo que se conoce como una \textit{topología} sobre $E$.

\medskip

Un \textit{espacio topológico} es un par ordenado $(X, \tau)$, donde $X$ es un
conjunto y $\tau$ es una colección de subconjuntos de $X$.
Para que $\tau$ sea una \textit{topología}, debe satisfacer los siguientes axiomas:

\begin{enumerate}
    \item $\emptyset \in \tau$ y $X \in \tau$.
    \item La unión de cualquier colección de conjuntos en $\tau$ también pertenece a $\tau$.
    \item La intersección de cualquier número \textit{finito} de conjuntos en $\tau$
    también pertenece a $\tau$.
\end{enumerate}

Los elementos de $\tau$ se llaman \textit{conjuntos abiertos}.

\medskip

Así, la topología no aparece aquí como una teoría independiente, sino como una
consecuencia estructural de la métrica inducida por el producto interno.
La geometría da lugar a la distancia, la distancia da lugar a los entornos,
y los entornos dan lugar a la estructura topológica.

\medskip

\textit{Interpretación geométrica.}

En $\mathbb{R}^n$ dotado de su producto interno estándar, las bolas abiertas
corresponden a las esferas familiares de la geometría euclídea.
Normas más generales deforman estas bolas en formas elipsoidales, reflejando
un escalado anisotrópico a lo largo de distintas direcciones del espacio.
En todos los casos, cada punto
\[
x = (x_1,\dots,x_n)
\]
es un vector del espacio vectorial ambiente, y los objetos geométricos
considerados son subconjuntos de ese espacio definidos por relaciones de
distancia entre sus vectores.

\medskip

Una vez que los objetos quedan representados como vectores de este modo, las
estructuras geométricas introducidas anteriormente ---normas, distancias, bolas
y elipsoides--- se convierten en herramientas para organizarlos y compararlos.

\medskip

Sus posiciones relativas, medidas mediante una métrica elegida, codifican
similitud.
Esta métrica no tiene por qué ser euclídea: en la práctica puede diseñarse
para reflejar invariancias propias del dominio o incluso aprenderse a partir
de datos.
Los entornos de significado se describen entonces mediante bolas abiertas,
y el aprendizaje puede interpretarse como el descubrimiento de una organización
geométrica y topológica subyacente a los datos.

\medskip

Una vez introducida una norma, ciertos conjuntos geométricos emergen como
consecuencias directas de la estructura y no requieren hipótesis adicionales.

\medskip

Sea $(E,\langle\cdot,\cdot\rangle)$ un espacio con producto interno y sea
$\|x\|=\sqrt{\langle x,x\rangle}$ la norma asociada.
Para un vector fijo $x_0\in E$ y un escalar $r>0$, la \emph{bola cerrada}
de radio $r$ centrada en $x_0$ se define como
\[
B_r(x_0)=\{\,x\in E \mid \|x-x_0\|\le r\,\}.
\]

La \emph{esfera} correspondiente es la frontera de este conjunto,
\[
S_r(x_0)=\{\,x\in E \mid \|x-x_0\|= r\,\}.
\]

\textit{Observación (Descripción coordenada en $\mathbb{R}^n$):}
Cuando $E=\mathbb{R}^n$ está dotado de su producto interno estándar y se
utilizan las coordenadas estándar $x=(x_1,\dots,x_n)$, las mismas definiciones
intrínsecas toman la familiar forma de «suma de cuadrados».
Para una bola centrada en el origen,
\[
B_r(0)=\left\{x\in\mathbb{R}^n\;\middle|\;x_1^2+\cdots+x_n^2\le r^2\right\},
\]
y su esfera frontera es
\[
S_r(0)=\left\{x\in\mathbb{R}^n\;\middle|\;x_1^2+\cdots+x_n^2=r^2\right\}.
\]
Esto es simplemente la expresión coordenada de
$\|x\|=\sqrt{x_1^2+\cdots+x_n^2}$.

\medskip

Estas definiciones son intrínsecas y no dependen de ninguna elección de
coordenadas o base.

\medskip

\textit{Observación (Simetría):}
En el caso euclídeo $E=\mathbb{R}^n$ con su producto interno estándar, la
bola es invariante bajo transformaciones ortogonales.
Si $Q\in O(n)$, entonces $\|Qx\|=\|x\|$, y por tanto
\[
Q\bigl(B_r(0)\bigr)=B_r(0).
\]
Más en general, toda isometría de un espacio con producto interno transforma
bolas en bolas.

\medskip

\textit{Observación (Secciones en $\mathbb{R}^n$):}
En $\mathbb{R}^n$, fijar la última coordenada $x_n=t$ muestra que la
sección transversal de una bola de radio $r$ por el hiperplano $x_n=t$ es
una bola $(n-1)$-dimensional de radio $\sqrt{r^2-t^2}$.
Esta observación contiene el contenido geométrico que subyace al método
habitual de «integrar por secciones» para calcular volúmenes $n$-dimensionales.

\medskip

\textit{Observación (Volumen de bolas euclídeas):}
En $\mathbb{R}^n$, el volumen $n$-dimensional de una bola euclídea escala como
una constante multiplicada por $r^n$:
\[
\operatorname{Vol}_n\bigl(B_r(0)\bigr)=\omega_n\,r^n,
\]
donde $\omega_n=\operatorname{Vol}_n\bigl(B_1(0)\bigr)$ es el volumen de la
bola unidad. Una fórmula cerrada es
\[
\omega_n=\frac{\pi^{n/2}}{\Gamma\!\left(\frac{n}{2}+1\right)}.
\]
En particular, $\omega_1=2$, $\omega_2=\pi$ y $\omega_3=\frac{4\pi}{3}$.

\medskip

\textit{Observación (Un efecto en alta dimensión):}
Aunque $B_1(0)\subset \mathbb{R}^n$ contiene siempre muchos puntos, su
volumen $n$-dimensional satisface $\omega_n\to 0$ cuando $n\to\infty$.
Esto ilustra un fenómeno básico de la geometría en alta dimensión: conjuntos
que parecen grandes en baja dimensión pueden ocupar una fracción despreciable
del volumen ambiente al aumentar la dimensión.

\medskip

Más en general, sea $B:E\times E\to \mathbb{R}$ una forma bilineal simétrica
definida positiva. El conjunto
\[
\mathcal{E}=\{\,x\in E \mid B(x-x_0,x-x_0)\le r^2\,\}
\]
se llama \emph{elipsoide}.
Así, un elipsoide es la bola unidad asociada a una forma bilineal, o
equivalentemente, a un producto interno modificado sobre $E$.

\medskip

Cuando la forma bilineal $B$ es simétrica y definida positiva, existe un
único operador lineal autoadjunto $A:E\to E$ tal que
\[
B(x,y)=\langle Ax,y\rangle
\]
para todo $x,y\in E$.

\medskip

El operador $A$ codifica la geometría del elipsoide.
Como $A$ es autoadjunto sobre un espacio con producto interno de dimensión
finita, admite una base ortonormal de vectores propios.
Si $\{e_1,\dots,e_n\}$ es tal base y $A e_i=\lambda_i e_i$ con
$\lambda_i>0$, entonces la desigualdad definidora
\[
B(x-x_0,x-x_0)\le r^2
\]
se convierte en
\[
\sum_{i=1}^n \lambda_i\, \xi_i^2 \le r^2,
\]
donde $\xi_i=\langle x-x_0,e_i\rangle$ son las coordenadas en esta
base propia.

% CORRECCIÓN 6: La observación de que la diagonalización «revela las direcciones
% invariantes ya presentes» aparecía dos veces en sucesión cercana. Se conserva
% la primera aparición (aquí, en la derivación del elipsoide) ya que concluye el
% argumento algebraico. La reiteración redundante que aparecía unas líneas más
% tarde, justo antes de la fórmula coordenada, ha sido eliminada.

\medskip

Así, los ejes principales del elipsoide son precisamente los vectores propios
de $A$, y los valores propios determinan el escalado a lo largo de esas
direcciones.
La diagonalización no crea los ejes; simplemente revela las direcciones
invariantes ya presentes en la estructura.

\medskip

En consecuencia, cuando se elige una base en la que la forma bilineal es
diagonal, esta definición invariante se reduce a una expresión coordenada
familiar.
En $\mathbb{R}^n$, se obtiene
\[
\mathcal{E}
=
\left\{
(x_1,\dots,x_n)\in\mathbb{R}^n
\;\middle|\;
\sum_{i=1}^n \frac{(x_i-a_i)^2}{h_i^2} \le r^2
\right\}.
\]

Escrita por componentes, esta desigualdad es
\[
\left(\frac{x_1-a_1}{h_1}\right)^2
+
\left(\frac{x_2-a_2}{h_2}\right)^2
+
\cdots
+
\left(\frac{x_n-a_n}{h_n}\right)^2
\le r^2.
\]

La ecuación depende de las coordenadas elegidas; el elipsoide en sí, no.

\medskip

\emph{Observación (Casos degenerados).}
La naturaleza geométrica del conjunto definido anteriormente depende de los
coeficientes $h_1,\dots,h_n$.
Si $h_1=\cdots=h_n$, el elipsoide se reduce a una esfera $n$-dimensional.
Si uno o más de los parámetros $h_i$ se anulan, la forma cuadrática
definidora se vuelve degenerada y el elipsoide colapsa en un objeto de
dimensión inferior: un disco, un segmento o, en el caso extremo, un punto.

\medskip

Desde el punto de vista estructural, estas degeneraciones corresponden a una
pérdida de rango en la forma cuadrática o bilineal asociada.
Así, los cambios en la estructura algebraica se traducen directamente en
cambios en la dimensión geométrica.

\medskip

\emph{Observación.}
A lo largo de esta discusión, el espacio ambiente $E=\mathbb{R}^n$ es un
espacio vectorial en el sentido del Capítulo I.
Cada elemento
\[
x=(x_1,\dots,x_n)
\]
es por tanto un vector de $E$.

\medskip

Sin embargo, los conjuntos definidos por estas expresiones son \emph{subconjuntos}
de $E$ seleccionados mediante restricciones cuadráticas.
No son espacios vectoriales en sí mismos, pues no son cerrados bajo la suma
ni bajo la multiplicación por escalares.

\bigskip

\textit{Resumen conceptual.}
Los espacios vectoriales proporcionan el escenario.
Las transformaciones lineales describen los cambios de estado admisibles.
La geometría aparece únicamente cuando se permite a los vectores interactuar
mediante una regla bilineal estructurada que produce escalares.
A partir de esta interacción emergen longitudes, ángulos, distancias y formas
geométricas.
La transición del álgebra a la geometría no es por tanto un salto, sino un
enriquecimiento controlado de la estructura.
El libro \cite{iribarren2} es una referencia excelente sobre la topología de
los espacios métricos.

\section*{Notas y Observaciones}

\textit{Perspectivas.}

En aplicaciones contemporáneas como el aprendizaje de representaciones,
los elementos de un conjunto (palabras, imágenes o señales) se proyectan en
espacios vectoriales de alta dimensión, donde las relaciones geométricas
codifican similitud e información estructural.
Estas construcciones proporcionan realizaciones concretas de ideas abstractas
del álgebra lineal.

\medskip

Los sistemas computacionales modernos ofrecen otro ejemplo ilustrativo.
Las plataformas dedicadas a la \emph{observabilidad} ---es decir, a la medición
sistemática del estado interno de un sistema en ejecución--- recopilan grandes
familias de cantidades numéricas como latencia, rendimiento, consumo de
memoria, tasas de error y utilización de recursos.

\medskip

En un instante de tiempo fijo, estas mediciones pueden ensamblarse en un vector
\[
x(t) = (x_1(t),\dots,x_n(t)) \in \mathbb{R}^n,
\]
cuyas coordenadas representan características observables del sistema.
A medida que evoluciona el tiempo, el sistema determina una trayectoria en este
espacio lineal ambiente, de modo que la monitorización y el análisis pueden
interpretarse geométricamente como el estudio de trayectorias en un espacio
vectorial de dimensión finita.

\medskip

Las consideraciones operacionales imponen típicamente cotas sobre cada observable,
expresadas mediante desigualdades de la forma
\[
a_i \le x_i \le b_i,
\]
que reflejan limitaciones físicas o rangos de operación aceptables.
Geométricamente, estas restricciones definen un producto cartesiano de
intervalos, es decir, un hiperrectángulo $n$-dimensional (o, tras
normalización, un hipercubo).
El comportamiento normal del sistema puede verse por tanto como confinado a
una región acotada del espacio ambiente, mientras que el comportamiento
anómalo corresponde a trayectorias que se alejan de dicha región.

\medskip

Conviene subrayar que esta construcción es un modelo matemático idealizado:
en los sistemas prácticos, la colección de observables puede variar con el
tiempo, por lo que la dimensión $n$ debe considerarse fija únicamente de manera
local o dentro de una representación elegida.
No obstante, el marco de los espacios lineales proporciona un lenguaje natural
para describir estados, restricciones y evolución.

\medskip

De este modo, los sistemas de observabilidad ilustran un tema recurrente del
álgebra lineal: un espacio ambiente plano sirve como escenario universal sobre
el cual los fenómenos significativos se caracterizan mediante subconjuntos
definidos geométricamente.
Las restricciones de norma dan lugar a esferas, los escalados anisotrópicos
producen elipsoides, y las cotas coordenadas definen hiperrectángulos.
A pesar de sus diferentes apariencias geométricas, todos ellos emergen como
subconjuntos estructurados embebidos en el mismo espacio lineal ambiente.

\bigskip
\hrule
\medskip

% CORRECCIÓN 7: La observación histórica sobre Lambert se ha trasladado a
% inmediatamente después del Interludio sobre Direcciones Invariantes en la
% sección de Notas, donde su conexión con los vectores propios y la preservación
% de ángulos es contextualmente inmediata. Anteriormente aparecía más adelante
% en las notas, separada de la discusión que la motivaba.

\textit{Observación histórica (Ángulos e invariancia).}

Antes de que los vectores propios se comprendieran como direcciones invariantes,
su detección solía ser indirecta.
En dos dimensiones, determinar si una transformación preserva una dirección se
reduce a detectar si preserva un ángulo.

\medskip

Este problema condujo de manera natural al cálculo de arcotangentes.
En el siglo XVIII, J.H.~Lambert introdujo la fracción continua
\[
\arctan(z)
=
\cfrac{z}{1 + \cfrac{z^2}{3 + \cfrac{4z^2}{5 + \cfrac{9z^2}{7 + \cdots}}}},
\]
en el marco de sus trabajos sobre mecánica orbital y fenómenos de rotación.

\medskip

Desde el punto de vista moderno, este esfuerzo analítico refleja una pregunta
geométrica: ¿rota una transformación lineal una dirección, o simplemente la
escala?
Los vectores propios proporcionan la respuesta estructural a esta pregunta,
haciendo del cálculo de ángulos algo auxiliar en lugar de fundamental.

\medskip
\hrule
\bigskip

\textit{Observación (Formas cuadráticas y el estadístico chi-cuadrado).}
Las expresiones de la forma que definen elipsoides aparecen de manera natural
en contextos que no son, a primera vista, geométricos.
Un ejemplo notable surge en la teoría de las pruebas de números aleatorios,
donde se encuentra la cantidad
\[
V
=
\frac{(Y_1 - np_1)^2}{np_1}
+
\frac{(Y_2 - np_2)^2}{np_2}
+
\cdots
+
\frac{(Y_n - np_n)^2}{np_n},
\]
conocida en estadística como el \emph{estadístico chi-cuadrado} asociado a
las cantidades observadas $Y_1,\dots,Y_n$ y los valores de referencia
$np_1,\dots,np_n$.

\medskip

Desde el presente punto de vista, esta expresión es simplemente una forma
cuadrática sobre $\mathbb{R}^n$.
Mide el cuadrado de la longitud del vector de desviación
\[
(Y_1 - np_1,\dots,Y_n - np_n)
\]
respecto a una forma bilineal simétrica definida positiva cuya matriz, en la
base estándar, es diagonal con entradas
$(np_1)^{-1},\dots,(np_n)^{-1}$.

\medskip

En consecuencia, los conjuntos definidos por desigualdades de la forma
\[
V \le c
\]
son elipsoides centrados en el punto $(np_1,\dots,np_n)$.
Los denominadores $np_i$ determinan el escalado anisotrópico a lo largo de
las direcciones coordenadas, exactamente como en el elipsoide general
\[
\sum_{i=1}^n \frac{(x_i-a_i)^2}{h_i^2} \le r^2.
\]

\medskip

Así, antes de imponer cualquier interpretación probabilística, la expresión
chi-cuadrado es ya un objeto geométrico: es la norma al cuadrado inducida
por un producto interno particular sobre $\mathbb{R}^n$.
La teoría estadística hace uso de esta geometría, pero no la crea.

\medskip

\textit{Observación (El espacio vectorial ambiente y las regiones restringidas).}
El propio espacio vectorial $E$ es el \emph{espacio ambiente} en el que viven
todos los vectores.
En coordenadas, cuando se elige una base, este espacio aparece como
$\mathbb{R}^n$, un entorno lineal plano en el que los vectores pueden sumarse
y escalarse libremente.

\medskip

Los objetos geométricos como bolas, esferas y elipsoides no reemplazan a este
espacio ambiente; son subconjuntos del mismo definidos por restricciones sobre
la norma o sobre una forma cuadrática.
Por ejemplo, la esfera
\[
S_r(x_0)=\{x\in E \mid \|x-x_0\|=r\}
\]
selecciona los vectores cuya distancia a $x_0$ es fija, mientras que un
elipsoide selecciona los vectores cuyas coordenadas satisfacen una desigualdad
cuadrática anisotrópica.

\medskip

En muchas aplicaciones modernas (el aprendizaje profundo en inteligencia
artificial, por ejemplo), los algoritmos operan sobre vectores en un espacio
ambiente de alta dimensión imponiendo tales restricciones durante el cálculo.
Los pasos de normalización pueden restringir los vectores a esferas, los
procedimientos de regularización pueden confinar los parámetros a regiones
elipsoidales, y los métodos de optimización buscan soluciones dentro de
entornos descritos por bolas.
El espacio vectorial ambiente proporciona el universo lineal en el que existen
los vectores, mientras que estos subconjuntos geométricos describen las regiones
seleccionadas por las reglas o restricciones del problema.